%\documentclass[preprint,12pt]{elsarticle}
\documentclass[final,3p,times,onecolumn]{elsarticle}
%\documentclass[preprint,5p,times]{elsarticle}
%\documentclass[preprint,12pt]{elsarticle}

\usepackage{amssymb}
\usepackage{makecell}
\usepackage{multirow}
\usepackage{booktabs}
\usepackage{multirow}
\usepackage{soul}
\usepackage{graphicx}
\usepackage{lscape}
\usepackage{subcaption}
\usepackage{booktabs}
\usepackage[T1]{fontenc}
\usepackage[utf8]{inputenc}
%\usepackage{lmodern}
\usepackage[english]{babel}
\usepackage[colorlinks]{hyperref}

\usepackage{longtable}
\usepackage{threeparttable} % For table notes
\usepackage{array}          % For p{} column widths
\usepackage{caption}     
\usepackage{amsmath}
\usepackage{float}
\usepackage[acronym,nomain,nonumberlist]{glossaries}
\loadglsentries{acronyms.tex}
\glsdisablehyper
% For better control of table captions
\captionsetup[table]{skip=10pt} % Adds some space between caption and table
%\input{glossary}
%\makeglossaries

%\newenvironment{framedglossary}%
%  {\begin{mdframed}}%
%  {\end{mdframed}}
%% The amsthm package provides extended theorem environments
%% \usepackage{amsthm}


%% The lineno packages adds line numbers. Start line numbering with
%% \begin{linenumbers}, end it with \end{linenumbers}. Or switch it on
%% for the whole article with \linenumbers.
%% \usepackage{lineno}

\journal{Energy Research \& Social Science }

\begin{document} 

\begin{frontmatter}

%% Title, authors and addresses

%% use the tnoteref command within \title for footnotes;
%% use the tnotetext command for theassociated footnote;
%% use the fnref command within \author or \address for footnotes;
%% use the fntext command for theassociated footnote;
%% use the corref command within \author for corresponding author footnotes;
%% use the cortext command for theassociated footnote;
%% use the ead command for the email address,
%% and the form \ead[url] for the home page:
%% \title{Title\tnoteref{label1}}
%% \tnotetext[label1]{}
%% \author{Name\corref{cor1}\fnref{label2}}
%% \ead{email address}
%% \ead[url]{home page}
%% \fntext[label2]{}
%% \cortext[cor1]{}
%% \affiliation{organization={},
%%             addressline={},
%%             city={},
%%             postcode={},
%%             state={},
%%             country={}}
%% \fntext[label3]{}

%Option 1 \title{Modeling Household Behavior in the Context of Solar Home System Electrification: A Case Study in Raqaypampa, Bolivia.}
%Option 2 \title{Energy, Culture, and Community: Understanding and Modeling Energy Behavior in Rural Raqaypampa, Bolivia}

% \title{Modeling Household Behavior in the Context of Solar Home System Electrification: A Case Study in Raqaypampa, Bolivia.}

\title{A Socio-Technical Framework for Modeling Household Energy Behavior in Remote Rural Electrification: A case in Raqaypampa, Bolivia}

%% use optional labels to link authors explicitly to addresses:
%% \author[label1,label2]{}
%% \affiliation[label1]{organization={},
%%             addressline={},
%%             city={},
%%             postcode={},
%%             state={},
%%             country={}}
%%
%% \affiliation[label2]{organization={},
%%             addressline={},
%%             city={},
%%             postcode={},
%%             state={},
%%             country={}}

\author[inst1,inst2]{Claudia Sanchez-Solis\corref{cor1}}
\ead{clsanchez@uliege.be}
\cortext[cor1]{Corresponding author}

\affiliation[inst1]{organization={Centro Universitario de Investigaciones en Energías, Universidad Mayor de San Simón},%Department and Organization
            %addressline={Address One}, 
            city={Cochabamba},
            %postcode={00000}, 
            %state={State One},
            country={Bolivia}}

\affiliation[inst2]{organization={Integrated \& Sustainable Energy Systems (ISES), Department of Mechanical and Aerospace Engineering, Université de Liège},%Department and Organization
            %addressline={Address Two}, 
            city={Liège},
            %postcode={22222}, 
            %state={State Two},
            country={Belgium}}

\affiliation[inst3]{organization={Instituto de Investigaciones en Ciencias Sociales, Universidad Mayor de San Simón},%Department and Organization
            %addressline={Address Two}, 
            city={Cochabamba},
            %postcode={22222}, 
            %state={State Two},
            country={Bolivia}}

\author[inst3]{Jaime Zambrana}
\author[inst1]{José Espinoza}
\author[inst1]{Sergio Balderrama}
\author[inst2]{Silvain Quoilin}

\begin{abstract}
%% Text of abstract
Accurate household demand estimation is decisive for the technical and financial viability of off-grid electrification, yet conventional models rely on urban-centric, income-driven assumptions that misrepresent the culturally embedded routines of remote rural and Indigenous communities. This paper develops and tests a mixed-methods pipeline that operationalizes qualitative socio-cultural knowledge to build a stochastic demand simulation. Drawing on a case study in the in Bolivia, and combining qualitative and quantitative data extracted under the Energy Cultures Framework analytical lens, the study extracts socio-technical rules governing daily energy use and consolidates them into four Energy Behavior Profiles. These profiles reveal an energy culture anchored in the agricultural cycle and are translated into behavior-specific parameters to build a stochastic demand model. Validation is performed per profile against measured data, then extended to a community-level comparison of two models: one composed of the socio-technically parametrized profiles, and a homogeneous one based on a single average profile. The first reproduces the shape of daily routines, the cross-profile ordering of day-to-day irregularity, and aggregate daily energy, properties a single community-average profile cannot capture by construction, yet that can directly affect system sizing. The paper argues that in data-scarce, culturally specific off-grid settings, qualitative understanding does not merely supplement quantitative modeling but structures it. The transferable contribution is the extraction-and-translation procedure itself, applicable to comparable contexts across the Global South.
\end{abstract}

% %%Graphical abstract
% \begin{graphicalabstract}
% %\includegraphics{grabs}
% \end{graphicalabstract}

% %%Research highlights
% \begin{highlights}
% \item Research highlight 1
% \item Research highlight 2
% \end{highlights}

\begin{keyword}
%% keywords here, in the form: keyword \sep keyword
Rural electrification \sep Solar Home Systems \sep Energy Culture \sep Developing countries \sep Energy Behavior \sep Energy Use

\end{keyword}

\end{frontmatter}

%% \linenumbers

%% main text
\section{Introduction}
\label{sec:introduction}

The global energy landscape is currently experiencing a transformation, driven by international commitments to climate change mitigation, sustainable development, and the eradication of severe energy poverty. Achieving universal access to clean, reliable, and affordable modern electricity remains one of the defining challenges of the twenty-first century. According to the 2025 Tracking SDG 7 Report, while global electricity access has reached 92\%, approximately 666 million people remain without power, and current trends indicate that 645 million will still lack access by 2030 \cite{iea_irena_unsd_wb_who_tracking_2023}. Furthermore, the \gls{iea} highlights that this deficit is heavily concentrated geographically and demographically: over 80\% of the global population lacking electricity resides in Sub-Saharan Africa, and 84\% live in remote, rural communities. Within this shifting macro-scenario, energy use around the world keeps evolving in time as a product of a broad range of factors that involve external conditions, technological advancements and personal decisions. Naturally, this evolution is heterogeneous and highly dependent on contexts. As a result, the energy consumption disparities persist, making decentralized renewable energy solutions—such as solar home systems and mini-grids—absolutely critical to reach isolated populations \cite{add cite}.

Accurate electricity demand estimation is a fundamental requirement for the sustainable planning, operation, and expansion of power infrastructure \cite{adinkrah_artificial_2025, poblete-cazenave_model_2021, guefano_forecast_2021}. It enables the execution of energy efficiency programs and maintains the balance between generation and consumption, a critical necessity given the difficulty of storing electricity at scale \cite{adinkrah_artificial_2025, kavousian_determinants_2013}. This precision becomes especially important in the context of small, decentralized off-grid networks like mini-grids and \gls{shs}, where forecasting is directly tied to a project's financial viability and technical survival \cite{lukuyu_purchasing_2023}. When demand is over-estimated, developers install oversized systems that drastically inflate upfront capital and maintenance costs, extending payback periods and risking financial collapse while unnecessarily increasing emissions \cite{lukuyu_purchasing_2023, blodgett_accuracy_2017, herraiz-canete_forecasting_2022, few_electricity_2022}. Conversely, under-estimating demand leads to undersized systems plagued by poor reliability, accelerated battery degradation, and the loss of critical services like healthcare \cite{lukuyu_purchasing_2023, blodgett_accuracy_2017}. Frequent outages erode community trust, force reliance on expensive diesel backups, and ultimately stunt organic economic growth \cite{herraiz-canete_forecasting_2022, few_electricity_2022}.

To avoid these pitfalls, off-grid operators must align community demand with the physical realities of generation and . In solar-powered microgrids, incentivizing daytime consumption when solar generation peaks can significantly reduce the need for massive battery storage, lowering the \gls{lcoe} by 25 percent or more \cite{lukuyu_purchasing_2023}. Achieving this requires a granular breakdown of demand to identify flexible loads—such as water pumping, agricultural milling, or cold storage—that can be shifted to match real-time generation \cite{kavousian_determinants_2013, few_electricity_2022}. Furthermore, it needs targeted Demand-Side Management strategies and time-of-use tariffs that are intentionally tailored to the micro-level financial behaviors of distinct customer segments, ensuring that newly stimulated productive uses complement rather than overwhelm the system's evening peaks \cite{lukuyu_purchasing_2023, lorenzoni_classification_2020}. In this context, knowing how energy is used is critical.

Traditionally, efforts to capture and simulate household energy consumption have relied heavily on a structural and economic lens, where income, socio-demographic makeup, and physical dwelling attributes serve as the primary drivers of demand \cite{pachauri_analysis_2004, Kowsari2011, poblete-cazenave_model_2021, richmond_drivers_2020}. While budget constraints, demographic footprints, and stark infrastructural divides dictate basic access and economies of scale \cite{jiang_rural_2020, kavousian_determinants_2013, zhao_survey_2019, azam_factors_2015, dominguez_estimating_2021}, reducing energy consumption to a purely technical or economic process critically misunderstands how human beings live. Extracting the level of granular, flexible demand required for off-grid survival is hard if planners only measure what and how much is being consumed; understanding the "why" behind energy consumption is absolutely critical. As energy literature increasingly recognizes, asking "What is energy for?" is the missing link in historical modeling efforts, which have traditionally analyzed abstract energy rather than human behavior \cite{Kowsari2011, shove_what_2014}. People do not consume raw electricity for its own sake; they consume "energy services" to accomplish deeply embedded social practices and everyday routines \cite{shove_what_2014, day_conceptualising_2016}.

This "why" reveals the cultural, moral, and social values that shape consumer choices, often overriding technical or economic logic \cite{Kowsari2011}. Households may reject efficient electric stoves because meals cooked over open flames are perceived as superior in taste and appropriateness \cite{kumar_non-universal_2020}; appliance acquisition may serve to claim modernity, mark social status, or strengthen community ties rather than meet a utility need \cite{tesfamichael_eager_2020}; and gathering biomass can function as a space for social bonding among women \cite{louw_determinants_2008}. Such practices are further differentiated by religion and ethnicity, which historically tie marginalized groups to particular fuels \cite{richmond_drivers_2020}. What top-down observers classify as "energy poverty" thus often reflects a localized, traditional lifestyle—energy behavior is a context-specific cultural construct \cite{kumar_non-universal_2020}. When planners treat rural users as purely rational economic actors, interventions fail: a new technology cannot replace a traditional one if it cannot perform the cultural "work" the household requires \cite{chattopadhyay_can_2024, shove_what_2014}.

Despite the deeply human reality of energy use, attempts to mathematically model and estimate energy demand face severe methodological hurdles that systematically minimize this human component. Accurate micro-level data is notoriously scarce, forcing modelers to rely on energy-use surveys that are highly vulnerable to self-reporting errors and "social-desirability bias" \cite{blodgett_accuracy_2017, m_heinrich_energy_2022}. While the integration of physical sensors in recent research has successfully exposed the massive biases in self-reported data, measured data alone cannot explain the complex behavioral dynamics driving those numbers; physical measurements must act as a complement to surveys, not a replacement \cite{rom_using_2020}. Lacking accurate behavioral data, energy simulation models frequently oversimplify users, treating them purely as rational economic actors seeking to maximize utility \cite{riva_long-term_2018}. This blind spot generates a significant "energy performance gap" between simulated predictions and actual real-world usage \cite{medojevic_integrated_2021}.

This discrepancy is exacerbated by the contextual misapplication of models developed in industrialized nations, which assume a linear progression up an income-based "energy ladder" and uniform power availability, failing to account for Global South realities like informal economies, suppressed demand, and fuel stacking—where exogenous factors such as unreliable grids compel households to maintain portfolios of traditional and modern fuels simultaneously \cite{stevanato_data-driven_2025, liu_impact_2022, Muhumuza2018}. Even bottom-up stochastic models, though mathematically sophisticated, struggle to define coincidence factors, either overestimating peak demand by assuming simultaneous appliance use or flattening the curve by spreading loads too evenly \cite{kuster_electrical_2017, herraiz-canete_forecasting_2022}. The problem runs deeper than income assumptions, however: the underlying Western, urban-centric paradigm treats households as utility-maximizing consumers, whereas critiques of this universalist ontology argue it misreads Indigenous and traditional rural communities \cite{shokrgozar_fine_2025}. Non-Western epistemologies such as \textit{Buen Vivir} instead frame well-being as relational—tied to community and environment rather than individual accumulation—so that what a model diagnoses as "energy poverty" or "suppressed demand" may reflect intentional sufficiency and cultural preservation \cite{gudynas_buen_2011}, while a "cosmologies of energy" approach ties energy use to cultural narratives, reciprocity, and traditional knowledge rather than mechanical work alone \cite{mazzone_indigenous_2022}. Models that ignore these ontological differences risk imposing inappropriate technical fixes and provoking technological rejection; advancing demand modeling in these settings therefore requires decentering the urban consumer baseline.

Consequently, standard energy demand models heavily rely on top-down, generalized assumptions that fail to capture the complex, highly localized realities of rural and indigenous communities. The \gls{ecf}, developed by Stephenson et al. \cite{stephenson_energy_2010, stephenson_energy_2015}, treats energy behavior not as a purely economic or technological choice but as a cultural phenomenon shaped by the dynamic interaction of three components: \textit{norms} (shared beliefs, expectations, and aspirations around energy use), \textit{material culture} (the technologies, buildings, and infrastructures involved), and \textit{practices} (the everyday activities that consume energy). These elements are embedded within a wider context of demographics, income, policy, and pricing that shapes and is shaped by the energy culture in turn. The framework is multi-scalar, having been applied from household transitions such as solar home system adoption \cite{khan_impacts_2020} to national decarbonization pathways \cite{stephenson_energy_2021}. For this work, its central value is diagnostic: the ECF is explicitly designed to characterize the heterogeneity that aggregate analyses obscure, clustering a population into distinct "energy cultures" or behavioral archetypes that share patterns of norms, practices, and material conditions \cite{stephenson_energy_2010, m_heinrich_energy_2022}. This makes it a natural foundation for identifying the behavioral profiles that structure energy demand in Raqaypampa.

In this context, critical methodological gap remains: conventional planning tools lack a systematic way to operationalize local cultural realities into the hard quantitative parameters required for accurate, inclusive electrification planning. To address this gap, the paper characterizes how socio-cultural factors shape energy behavior in Indigenous communities and operationalizes these insights for bottom-up demand modeling. Drawing on a detailed case study of the Indigenous Territory of Raqaypampa, Bolivia, the study introduces a novel mixed-methods pipeline that translates qualitative insights into parameters to build a demand simulation model, replacing generalized assumptions with localized empirical reality. 

The key contributions of this paper are the following:

\begin{itemize}
    \item Applying the Energy Cultures Framework, the paper describes energy behavior in an isolated Indigenous Quechua community and the socio-cultural, economic, and material factors shaping daily energy practices.
    \item The paper sets out a procedure that integrates qualitative and quantitative data to produce parameters to build a stochastic bottom-up demand model, capturing behavioral particularities and community heterogeneity.
\end{itemize}

\section{Methods}
\label{sec:methodology} 
Mixed methods combine qualitative and quantitative data through rigorous collection, analysis and through deliberate procedures for integrating the two \cite{creswell_research_nodate}. The design has been applied across energy research, from technology adoption \cite{university_of_augsburg_adoption_2019, palm_peer_2017} to the study of everyday energy practices \cite{rosenberg_evidence_2019, guerra-santin_mixed_2016, wolff_tackling_2017}. In this sense, this study adopts a mixed-methods approach, analyzing several data streams independently before integrating them to derive the quantitative parameters that feed a stochastic demand model.  Figure~\ref{fig:methods} presents the methods flowchart. The following sections detail the process: the study area, the multi-stream data collection, the integration and the modeling procedure through which the collected data becomes a validated demand model.

\begin{figure}
    \centering
    \includegraphics[width=0.9\linewidth]{manuscript/Figures/Figure_1.png}
    \caption{Methods flowchart}
    \label{fig:methods}
\end{figure}

\subsection{Study Area: An Indigenous Originary Peasant Territory}
\label{sec:case_study}

This study focuses on an electrification initiative implemented across three communities located within the \gls{tioc}\footnote{From the Spanish \emph{Territorio Ind\'igena Originario Campesino}.} of Raqaypampa in Bolivia, South America. Following collaborative, consensus-based dialogues with community leaders, 100 systems were deployed featuring a solar panel, a battery, and LED lighting accessories (see Table \ref{tab:shs_characteristics}). Each household owning a \gls{shs} is marked by a yellow dot on the map in Figure~\ref{fig:raq_map}.

\begin{figure}
    \centering
    \includegraphics[width=0.5\linewidth]{manuscript/Figures/Figure_2.png}
    \caption{Indigenous Originary Peasant Territory of Raqaypampa}
    \label{fig:raq_map}
\end{figure}

\begin{table}[]
\centering
\begin{tabular}{ll}
\hline
Component  & Characteristics                                                          \\ \hline
Solar PV   & \begin{tabular}[c]{@{}l@{}}Polycrystalline\\ 20Wp, 12V \end{tabular}       \\ \cline{1-2}
Battery    & \begin{tabular}[c]{@{}l@{}}LiFePO4\\ 12.8V \\ 7000mAh, 89Wh\end{tabular} \\ \cline{1-2}
\multirow{5}{*}{Accessories} & 1 LED light 2W, 200 lm SFL                                               \\
           & 1 LED light 3W, 200 lm SFL                                               \\
           & Phone charging set                                                       \\
           & Metal structure                                                          \\
           & Rechargeable radio                                                       \\ \hline
\end{tabular}
\caption{Characteristics of the Solar Home System acquired by the families in Raqaypampa.}
\label{tab:shs_characteristics}
\end{table}
Officially recognized as the first \gls{tioc} in the country on July 12, 2016, Raqaypampa marked a milestone in indigenous self-determination. Its Autonomy Statute was approved by referendum in November 2016, with over 91\% voter support, followed by the first elections of autonomous authorities in June 2017. Grounded in ancestral values and the principle of Sumaq Kawsay (Living Well), the territory is governed through holistic, culturally rooted practices aimed at managing the land in harmony with \textit{Pachamama} (Mother Earth). Today, Raqaypampa serves as a national reference for community-led governance. The territory encompasses 43 communities across approximately 556 square kilometers, with elevations ranging from 1,600 to 3,300 meters above sea level (See Figure \ref{fig:raq_map}) \cite{oep_estatuto_2016, gaioc_del_territorio_de_raqaypampa_plan_2018}. 

Between 2001 and 2012, the territory experienced a significant population decline, dropping from approximately 10,230 to 7,340 inhabitants. This decrease was largely attributed to permanent migration driven by economic hardship and the absence of basic services. However, according to the most recent national census conducted in 2024, the population has increased slightly, reaching 8,007 inhabitants \cite{CENSUS2024}. 

Household sizes in the targeted communities range from 1 to 12 members. Families are predominantly organized as nuclear, extended, or large extended households, with 75\% led by male heads of household. Dwellings are typically constructed from mud bricks with earthen floors and roofs of clay tile, corrugated metal, or straw. Living spaces are generally shared, with separate auxiliary structures frequently built for cooking with firewood and storing food (known locally as \textit{trofas}). An example of this structural distribution is shown in Figure \ref{fig:typical_household}.

The local economy relies heavily on rain-fed, subsistence agriculture and livestock farming, with occasional surplus sold in local markets. Livelihoods remain highly precarious; at least 60\% of households report earning the equivalent to less than USD 35 per month or having no direct monetary income. The communities lack access to potable water infrastructure, relying instead on natural springs and rivers that are highly vulnerable to drought. Diets are largely based on locally grown maize, potatoes, and vegetables, supplemented by rice.

Because the communities lack irrigation, survival and economic stability are strictly tethered to the rainy season. This establishes a rigid annual baseline calendar: soil preparation occurs in August–September, sowing in October–November, crop maintenance throughout the rains (October–March), and harvesting in May–June.

\begin{figure}
    \centering
    \includegraphics[width=0.5\linewidth]{manuscript/Figures/Figure_3.png}
    \caption{Typical household in Raqaypampa}
    \label{fig:typical_household}
\end{figure}



\subsection{Multi-Stream Data Collection}
\label{sec:data_collection_analysis} 

The framework of the Energy Cultures proposed by Stephenson et al. \cite{stephenson_energy_2010} is used to guide data gathering tools design. Table \ref{tab:variables} in the Supplementary Material shows the operationalization of each of its dimension. The three main data streams are described as follows:

\begin{itemize}
    \item \textbf{Stream A: Qualitative data}. 
    Comprises the study's qualitative material: semi-structured interviews with household heads and key informants, capturing daily habits, gender roles, energy-sharing norms, and seasonal traditions, alongside dated field memos kept by the lead researcher. Each memo is tied to a specific household and records structural and demographic observations, forming a distinct, attributable evidence tier that corroborates the interview material.
    \item \textbf{Stream B: Socioeconomic and Self- reported energy practices surveys}. 
    Comprises the study's structured quantitative material, capturing the demographic boundaries and self-reported behaviors of the community through two survey instruments. A baseline socioeconomic survey records demographics, livelihoods, and structural constraints, while a follow-up practices survey systematically maps self-reported energy behaviors across the community.
    \item \textbf{Stream C: Empirical Load Data}. 
    Comprises the study's empirical load data, capturing the objective physical reality of energy consumption through continuous data loggers installed in a subset of \gls{shs} households. These sensors record real-time, high-resolution electricity consumption, from which the true baseline wattage and the precise timing of appliance use are extracted. This measured data serves as the validation source against which the modeled load curves are assessed.
\end{itemize}

For this study, the data collection campaign—comprising the socioeconomic and practice surveys (Streams~A and~B), semi-structured interviews and field memos (Stream~A), focus group discussions, and datalogger deployment (Stream~C)—was carried out between 2023 and 2026, as detailed in Table~\ref{tab:data_collection}.

Qualitative data (Stream A) is coded in the open-source software QualCoder using a hybrid deductive--inductive approach: the three core dimensions of the Energy Cultures framework seeded a foundational codebook, while context-specific sub-codes emerged inductively from the transcripts. The coded data then passes through a three-step analytical pipeline. Code-frequency reports within the Practices category first isolate the high-frequency energy events that anchor the community's physical load profile. Co-occurrence queries run with QualCoder's matrix tools then map these temporal anchors against demographic vectors and cultural norms, identifying which household profiles drove specific peak-demand events and why. Finally, these intersections were translated into analytical memos following a standardized formula:
\[
[\text{Norm/Demographic}] \;\text{ drives }\; [\text{Practice}]
\;\text{ under }\; [\text{Material Conditions}].
\]
It expresses each socio-cultural driver as an explicit logic rule. The full coding procedure, codebook, and matrix queries are provided in the Supplementary Material S1. Through this process, the thematic synthesis yields a discrete set of rigid social rules governing daily energy use, which forms the base for Energy Behavior Profile characterization. Stream B data is used in the integration part as described in the following sections.

Empirical load data (Stream~C) from the physical sensor network provide high-resolution time-series records of energy demand. The raw datasets are first cleaned using Python libraries to ensure integrity and standardize temporal intervals, after which the measured load curves are analyzed quantitatively and used for validation.

\begin{table}[ht]
\centering
\small
\setlength{\tabcolsep}{4pt}
\renewcommand{\arraystretch}{1.1}
\caption{Overview of data collection methods applied in the Raqaypampa case study. All households are SHS beneficiaries; instruments and interview guides are in Appendix~A.}
\label{tab:data_collection}
\begin{tabular}{@{}>{\raggedright\arraybackslash}p{2.7cm}
                   >{\raggedright\arraybackslash}p{2.7cm}
                   >{\raggedright\arraybackslash}p{2.6cm}
                   >{\raggedright\arraybackslash}p{5.0cm}@{}}
\toprule
\textbf{Method} & \textbf{Sample / participants} & \textbf{Timing} & \textbf{Purpose} \\ 
\midrule
Structured survey I & 100 household heads & Feb 2023 (baseline) & Socioeconomic data, household characteristics, energy access. \\
Structured survey II & 65 household members & Nov 2023--Dec 2025 & Self-reported energy practices. \\
Semi-structured interviews & 50 (heads and key informants) & Mar 2023--Feb 2026 & Energy practices, cultural norms, perceptions of sufficiency. \\
Focus group discussions & 2 groups (10--15 each) & Apr 2023 & Validate survey findings; gender roles, energy-sharing, aspirations. \\
Electricity data loggers & 14 SHS households & Continuous, Apr 2023--Mar 2026 & Real-time generation/consumption for demand-model validation. \\ 
\bottomrule
\end{tabular}
\end{table}

\subsection{Data integration: Energy Behavior Profile characterization}

\glspl{ebp} integrate the qualitative socio-cultural rules (Stream~A) with demographic survey variables (Stream~B). During thematic coding, subsets of the extracted socio-technical rules were observed to co-occur consistently within the same types of households, indicating a latent typology in which distinct rule-bundles govern distinct segments of the population rather than applying uniformly. These emergent clusters were confirmed through the relational matrix analysis of Section~\ref{sec:data_collection_analysis} and form the origin of the profiles.

The variables characterizing each profile are read directly from the rule formulations rather than selected through statistical screening: because each rule is framed around specific actors and social roles, the actors it names are matched to their corresponding survey fields to yield a small set of structural proxies. These proxies extend the typology from the interviewed households to the full surveyed population, operationalized as a fixed, priority-ordered decision tree evaluated in strict order and stopping at the first match, so that higher-priority variables capture the structural factors dominating a household's routine and lower-priority variables resolve residual ambiguity. The tree is applied only to households with complete coverage of every required field across the relevant survey waves---a deliberate scope constraint, not incidental sample loss---and missing fields are resolved by a sensitivity check rather than a silent default: every plausible value is substituted and the tree re-run, assigning a profile only where all combinations agree and leaving the household explicitly unclassified otherwise. The variable ordering, decision logic, and resulting counts are reported in Section~\ref{sec:ebp_results}.

The classification is validated in two independent ways: a qualitative concordance check screens each classified household's interview and memo text against the thematic markers of its assigned profile, flagging disagreements for manual review without editing the classification; and a quantitative correlation analysis tests discriminant validity, confirming that the groups are behaviorally distinct along the dimensions their governing rules predict (Section~\ref{sec:ebp_results}). Because the classification criteria are entirely survey-derived, every household within scope is assigned to exactly one profile, while only the interviewed subset supplies the qualitative material used in the parametrization that follows (Section~\ref{sec:va_parametrization}). The full procedure, per-profile counts, and survey-derived versus interviewed subpopulations are reported in Supplementary Material~S1.

\subsection{Simulation of the energy demand}
\label{sec:ramp}

Once the data streams are integrated and households are classified into user types according to the socio-technical rules that shape their energy practices, the demand model is built. The load curves are generated with RAMP, an open-source bottom-up stochastic model, and validated against the measured data of Stream C. The following sections describe the tool, the parametrization protocol and the validation strategy.

\subsubsection{RAMP: Multi-Energy Load Profile Generation}
\label{sec:ramp}
The simulations are implemented in RAMP, an open-source, bottom-up stochastic model for generating high-resolution multi-energy load profiles in contexts where metered data is scarce or absent \cite{lombardi_ramp_2024}. Rather than reproducing a fixed demand curve, RAMP builds each daily profile probabilistically from the ground up: every appliance is defined by its power, the number of daily functioning windows, the total time it is switched on, and the time windows within which that use may occur, and the model samples usage stochastically within these bounds to produce a distinct load profile for each simulated day. This appliance-level, probabilistic construction makes RAMP well suited to the present study, since it accepts exactly the behaviorally derived parameters—time windows, activation probabilities, and usage durations—that the Virtual Appliance protocol produces, and its day-to-day stochasticity reproduces the natural variability of human routines rather than collapsing them to an average.

\subsubsection{Virtual Appliance parametrization}
\label{sec:va_parametrization}

Integrating the qualitative rules of the Energy Cultures Framework into a stochastic engine requires a shift in how appliances are represented. Conventional models standardize hardware, assigning a single generalized distribution to a generic device (e.g., a "3W LED bulb") spread across the day. Grounded in Social Practice Theory, this study instead treats energy use as the performance of social practices \cite{shove_what_2014} and introduces the \gls{va} to operationalize this in a bottom-up engine: not a physical device but a mathematically bounded human routine, in which a single bulb is disaggregated into multiple behavior-specific loads—for example a "Morning Agricultural Prep Light" and an "Evening Domestic Cooking Light"—each with its own empirically derived time window and activation probability, dictated by the norms and material constraints of a household's profile.

Each \gls{va}'s time window is set by first reconstructing practice timing qualitatively as an Anthropological Window—a culturally bound segment of the household's day (e.g., the evening domestic block) built from interview narrative—and then triangulating it against independent survey time-use fields and physical anchors (sunrise/sunset). Because narrative and survey evidence are partially independent, their agreement is itself checkable and is reported as a triangulation consistency finding where the data permit.

The translation from qualitative account to quantitative parameter follows a systematic protocol in which every value traces to one of five declared derivation sources (Table~\ref{tab:derivation_sources}): two direct technical inputs, the triangulated Anthropological Window, and two crosswalks that sort respondent language into quantitative bins—one mapping frequency language to the activation probability \texttt{occasional\_use}, the other mapping a household's account of a practice into a rigidity bin (Strict / Flexible / Chaos).

\begin{table}[htbp]
\centering
\small
\setlength{\tabcolsep}{5pt}
\renewcommand{\arraystretch}{1.2}
\caption{Derivation sources for Virtual Appliance parameters.}
\label{tab:derivation_sources}
\begin{tabular}{@{}>{\raggedright\arraybackslash}p{2.0cm}
                   >{\raggedright\arraybackslash}p{6.0cm}
                   >{\raggedright\arraybackslash}p{4.6cm}@{}}
\toprule
\textbf{Source} & \textbf{Basis} & \textbf{Parameters} \\
\midrule
{[SPEC]} & Hardware specification (non-behavioral) & \texttt{power}, \texttt{number} \\
{[WINDOW]} & Anthropological Window: interview $\times$ survey triangulation, physically anchored & \texttt{window\_1}, \texttt{window\_2} \\
{[FREQ-XW]} & Frequency crosswalk: respondent language $\rightarrow$ activation probability & \texttt{occasional\_use} \\
{[RIG-XW]} & Rigidity crosswalk: account $\rightarrow$ bin (Strict / Flexible / Chaos) & \texttt{time\_fraction\_\-random\_variability}, \texttt{random\_var\_w}, \texttt{func\_cycle}, \texttt{func\_time} \\
{[OCC]} & Household occupancy: structural absence from interview narrative, suppressing all VAs at once & \texttt{prob\_home} \\
{[DECLARED DEFAULT]} & Stated, reasoned value where no source speaks to the parameter & \texttt{power\_var} \\
\bottomrule
\end{tabular}
\end{table}

The rigidity crosswalk illustrates the paper's core argument. A single qualitative judgment of how rigid a practice is—anchored to an external constraint (Strict), shifting but describable (Flexible), or unpredictable from the respondent's own account (Chaos)—resolves three separate model parameters: two directly, and \texttt{func\_cycle} via a second-order fraction of \texttt{func\_time}. One qualitative judgment thus does the work of three independently justified quantitative estimates, so qualitative insight does not merely supplement the model but provides its framework. Each value reflects behavior as reported by respondents rather than an independently validated ground truth, and traces to a dated interview anchor. The complete crosswalk tables, worked derivations with anchor quotes, and the protocol's revision history are provided in Supplementary Material~S1.

\subsubsection{Validation strategy}
\label{sec:methods_sim}

To isolate the value of integrating qualitative data into quantitative demand modeling, a comparative validation strategy evaluates whether socio-technically derived parameters capture behavioral detail that improves the model. Two steps are followed: a per-profile validation against measured data, and a community-scale comparison that isolates the effect of modeling household diversity.

Load curves are simulated with RAMP for each of the four Energy Behavior Profiles, using the Virtual Appliance parameters of Section~\ref{sec:va_parametrization}, and compared against a measured baseline built by averaging the available Stream~C datalogger records of each profile's classified households. Because decentralized household load curves are highly volatile and driven by human activity, point-by-point magnitude alone is an inadequate benchmark, so the comparison relies on a Dual-Tier Validation Strategy (Table~\ref{tab:indicator_selection}). Tier~1 assesses daily chronological alignment: \gls{rmse} measures step-by-step deviation but heavily penalizes minor temporal shifts in human routines, so it is paired with \gls{lcss}, which rewards shape and sequence alignment while tolerating behavioral elasticity, and with \gls{mpdada}, which isolates peak-demand events. Tier~2 complements the analysis validating the emergent structural properties a chronological error cannot see---the sustained near-peak level that sets converter ratings (P95 power), total daily energy that sizes PV and battery capacity, the day-to-day scatter of the peak instant that signals a household's rigidity (peak-time spread), and the day-to-day structural irregularity of routines represented by the Chaos Index or \gls{MRSD}. Complementing these, the \gls{LDC}---demand sorted in descending order against the fraction of the day at or above each level, discarding timing entirely---is read graphically to assess how faithfully each model reproduces the distribution of demand levels a household sustains. Full indicator definitions, formulas, and sources are given in Table~\ref{tab:indicator_selection}.

\begin{table}[htbp]
\centering
\caption{Taxonomy of the load-curve indicators selected for the Dual-Tier Validation strategy, with their socio-technical meaning.}
\label{tab:indicator_selection}
\small
\renewcommand{\arraystretch}{1.3}
\begin{tabular}{>{\raggedright\arraybackslash}p{2.5cm} >{\raggedright\arraybackslash}p{3cm} >{\raggedright\arraybackslash}p{5.5cm} >{\raggedright\arraybackslash}p{3cm}}
\toprule
\textbf{Validation Tier} & \textbf{Indicator} & \textbf{Socio-Technical / Behavioral Meaning} & \textbf{Reference(s)} \\ 
\midrule
\textbf{Tier 1 Validation} \newline \textit{(Time-Series)} & Root Mean Square Error (RMSE) & Evaluates the strict mathematical magnitude of step-by-step simulation error; reported normalized by mean power (nRMSE) and benchmarked against each household's own variability for cross-profile comparability. & Köhler et al., 2022; Nti et al., 2020; Adeoye et al., 2019; Widén et al., 2009 \\
& Longest Common Subsequence (LCSS) & Assesses chronological sequence alignment; accommodates human behavioral elasticity and minor time shifts. & Zhang et al., 2022; Köhler et al., 2022 \\
& Mean Percentage Deviation Against Daily Average (MPDADA) & Isolates peak demand periods; explicitly measures the model's capacity to overcome the flattening ``Pancake Effect.'' & Ihsane et al., 2018 \\
\midrule
\textbf{Tier 2 Validation} \newline \textit{(Structural Profiling)} & P95 Power Error & Captures the sustained near-peak demand level (exceeded only 5\% of the day); read from the LDC. & \textit{[add ref]} \\
& Daily Energy Error & Validates total daily energy consumed (equivalently, mean power). & \textit{[add ref]} \\
& Peak-Time Spread & Day-to-day scatter (standard deviation, in hours) of the daily peak instant; a rigidity signal, testing reproduction of timing dispersion rather than average peak placement. & [add ref] \\
& Mean Relative Standard Deviation (MRSD) & Validates day-to-day structural irregularity, capturing nomadic seasonal migration or informal hardware modifications. & Kaur et al., 2022 \\
\bottomrule
\end{tabular}
\end{table}

To complement the per-profile validation, the analysis is extended to the community scale, since sizing decisions act on the aggregate load a system serves rather than on any single household. Two constructions of community demand are contrasted:

\begin{itemize}
\item \textbf{Heterogeneous community (Socio-Technical)}: the four Energy Behavior Profiles aggregated in their classified proportions, each parametrized from its Virtual Appliances.
\item \textbf{Homogeneous community (Survey-based)}: a single average profile, informed only by community-wide survey averages, replicated across all households.
\end{itemize}

Both are simulated over one year at the classified population size. This is a comparison rather than a validation: no aggregated measured data exists at the community level, so the heterogeneous aggregate---the community-scale composition of the per-profile models already validated above---serves as the empirically anchored reference against which the homogeneous model is assessed. The comparison isolates one effect in particular: the partial cancellation of non-coincident peaks, invisible at the single-profile scale and absent by construction from a homogeneous model. It is quantified by the coincidence factor---the aggregate peak divided by the sum of individual household peaks---which measures how much diversity in peak timing lowers the demand a shared system must meet. The aggregate is further characterized with the same indicators used per profile applied to the summed community load, together with a graphical reading of the community-level \gls{ldc}.

\section{Results}
\label{sec:results}

Results are presented in three parts, moving from qualitative pattern to quantitative validation. Section~\ref{sec:energy-culture} identifies the socio-technical rules governing household energy behavior and the demographic variables that distinguish them, converging on four \glspl{ebp}. Section~\ref{sec:va_parametrization} translates qualitative and quantitative data into the \gls{va} parameters used for the RAMP simulation, disaggregated physical hardware into behaviorally distinct loads for the four profiles. Section~\ref{sec:energy_demand_model} presents the results of simulating load curves for the individual profiles and hole communities, then evaluates whether this socio-technical parametrization improves on a simplified baseline against measured data, across both aggregate accuracy and structural load-shape metrics.

\subsection{The Energy Culture of Raqaypampa}
\label{sec:energy-culture}

The integration of qualitative interviews (Stream A) and structured household surveys (Stream B) reveals that electricity demand in Raqaypampa is not randomly distributed, nor is it driven purely by appliance ownership. Instead, energy behavior is governed by a matrix of cultural norms, demographic shifts, and material constraints, and this governance is legible in how households describe their own routines. Through relationship-based thematic analysis of code co-occurrences, 17 distinct socio-technical rules are extracted. These rules act as the governing logic for daily energy consumption and are organized into four operational themes as shown in Table \ref{tab:socio_technical_rules}.

\begin{table}[]
\centering
\footnotesize
\setlength{\tabcolsep}{4pt}
\renewcommand{\arraystretch}{1.2}
\begin{tabular}{@{}clp{0.62\textwidth}@{}}
\toprule
\textbf{\#} & \textbf{Rule} & \textbf{Governing logic (driver $\rightarrow$ practice under material condition)} \\
\midrule
\multicolumn{3}{@{}l}{\textit{Daily Baselines --- Predictable Anchors}} \\
\midrule
1 & Educational Anchor & School-aged children $\rightarrow$ protected evening lighting block (18:00--22:00) on one shared LED \\
2 & Agricultural Dictate & Rain-fed farming calendar $\rightarrow$ pre-dawn activity (04:00--05:00) under early-morning lighting \\
3 & Gendered Anchor of Domestic Operations & Adult women's food prep $\rightarrow$ sets onset of evening peak under kitchen lighting \\
4 & Connectivity Lifeline & Out-migration of children $\rightarrow$ phone charging as non-discretionary need under limited battery \\
5 & Digital Delegation & Presence of youth/school-age member $\rightarrow$ charging cadence tracks that member, not nominal owner \\
\midrule
\multicolumn{3}{@{}l}{\textit{System Constraints \& Coping Mechanisms}} \\
\midrule
6 & Weather-Driven Load Shedding & Cloudy days $\rightarrow$ ration lighting, shift charging to daylight under supply scarcity \\
7 & Safety Baseline & Documented nocturnal harm $\rightarrow$ $\geq$1 LED left on until dawn, raising overnight DoD \\
8 & Infrastructure Mobility & Seasonal field labor $\rightarrow$ panel dismantled and hauled (``trajinar'') to camps; multi-day load drop-offs \\
9 & Nature as a Threat & Fear of lightning $\rightarrow$ maintenance suspended in rainy season under exposed siting \\
\midrule
\multicolumn{3}{@{}l}{\textit{Structural Anomalies}} \\
\midrule
10 & Aspiration Gap \& Dual-Home Displacement & Access to 2nd grid-connected residence $\rightarrow$ heavy loads relocate; SHS home zoned for survival (flat curve) \\
11 & Protective Gatekeeping & Fragile, irreplaceable asset $\rightarrow$ sharing refused despite reciprocity norms \\
12 & Community Override & Major communal events (e.g.\ \textit{velorio}) $\rightarrow$ conservation suspended; rare high-impact spikes \\
13 & Hardware Resignation & Component failure + absent repair/supply chains $\rightarrow$ permanent load reduction, not repair \\
\midrule
\multicolumn{3}{@{}l}{\textit{Productive Labor \& Multi-Energy Realities}} \\
\midrule
14 & Productive Lighting Veto & LED inadequate for detailed handwork $\rightarrow$ no evening weaving load; weaving stays in daylight \\
15 & Fuel Stacking & Wet firewood / roadless terrain $\rightarrow$ default biomass, switch to hoarded LPG when wet \\
16 & Gratitude/Scarcity Baseline & Memory of pre-electrification (\textit{mechero} smoke) $\rightarrow$ conservative, finite-good use of light \\
17 & Sedentary Anchor & Aging / physical limitation $\rightarrow$ exempt from seasonal migration; continuous year-round baseline \\
\bottomrule
\end{tabular}
\caption{The seventeen socio-technical rules governing daily energy use in Raqaypampa, grouped by operational theme. Full anchor quotes for each rule are provided in Supplementary Material S1.}
\label{tab:socio_technical_rules}
\end{table}


\subsubsection{Energy Behavior Profiles (EBP)}
\label{sec:ebp_results}

Before consulting any survey variable, four profiles emerge from thematic coding of the interviews: specific rule-bundles are observed to co-occur consistently within the same types of households, rather than distributing uniformly across the sample (See Table \ref{tab:ebp-summary}). The Educational Anchor (Rule 1), Agricultural Dictate (Rule 2), Gendered Anchor of Domestic Operations (Rule 3), and Digital Delegation (Rule 5) consistently co-occur among working-age nuclear households with school-aged children, motivating the \gls{p1}. The Connectivity Lifeline (Rule 4) and Sedentary Anchor (Rule 17) co-occur among elder-headed households with no co-resident working-age adults or young children, inspiring the \gls{p2}. The extended-cooking variant of the Gendered Anchor (Rule 3), Aspiration Gap and Dual-Home Displacement (Rule 10), Community Override (Rule 12) and Safety Baseline (Rule 7) consistently co-occur among multi-generational households under one roof, driving the \gls{p3}. Infrastructure Mobility (Rule 8), aspiration Gap and Dual-Home Displacement (Rule 10) and Hardware Resignation (Rule 13) consistently co-occur among households whose occupation structurally decoupled them from the settlement's routine, motivating the \gls{p4}.

\begin{table}[H]
\centering
\small
\setlength{\tabcolsep}{4pt}
\caption{Summary of Energy Behavior Profiles (EBPs) and their expected stochastic load signatures. Governing rules refer to Table~\ref{tab:socio_technical_rules}.}
\label{tab:ebp-summary}
\begin{tabular}{@{}>{\raggedright\arraybackslash}p{2.6cm}
                   >{\raggedright\arraybackslash}p{4.5cm}
                   >{\raggedright\arraybackslash}c
                   >{\raggedright\arraybackslash}p{4.5cm}@{}}
\toprule
\textbf{Profile (EBP)} & \textbf{Key Demographics} & \textbf{Rules} & \textbf{Expected Load Signature} \\
\midrule
P1: Educational \& Agricultural Core &
Nuclear or stable extended families (no migration, SHS kept at home) with school-aged children &
1, 2, 3, 5 &
Sharp early-morning spike; midday valley; heavy, rigid evening peak; high daily phone charging. \\
\addlinespace
P2: Isolated Elderly &
Unipersonal or nuclear-couple households; elders aging in place, no young children &
4, 7, 17 &
Later, smoother morning wake-up; flat daytime radio load; persistent 1--3\,W overnight safety lighting. \\
\addlinespace
P3: Extended \& Multi-Tasking Hub &
Extended-family households failing Profile~1's stability clause (active migration or portable SHS) &
3, 10, 12 &
Longest evening lighting (extended cooking); extreme daytime USB competition; highest early-depletion risk. \\
\addlinespace
P4: System Breakers &
Non-agricultural household head (overrides family structure), plus residual households &
8, 10, 13 &
Highly stochastic: load drops to zero for days/weeks during migrations, or false baselines from parallel infrastructure. \\
\bottomrule
\end{tabular}
\end{table}


To operationalize this qualitative typology for the full surveyed population, each profile is mapped to a specific structural proxy, read from the actors and material conditions the rules themselves name: \texttt{family\_type} (naming the household-structure actors of Rule 3 and the Structural Anomalies rules, Rules 10 and 12), \texttt{occupation} (naming the mobile and adaptive actors of Infrastructure Mobility, Rule 8), \texttt{children\_in\_school} (naming the actors of the Educational Anchor, Rule 1), \texttt{migration}, and \texttt{portability\_shs} (jointly naming the stability conditions distinguishing a settled extended household from a mobile one).

The five splitter variables identified above are applied as a fixed, priority-ordered decision tree, evaluated top to bottom with the first matching rule assigning the profile. The tree is applied to the $N = 65$ households with complete coverage across both survey waves; this scope constraint follows from \texttt{portability\_shs} being collected only in the social-practices instrument, and is a deliberate design decision rather than incidental sample loss. Table~\ref{tab:tree} states the routing logic and the resulting distribution.

\begin{table}[ht]
\centering
\caption{Household routing logic and resulting profile distribution ($N=65$). Steps are evaluated in order; the first match assigns the profile.}
\label{tab:tree}
\begin{tabular}{clcc}
\toprule
Step & Test on the household & Profile & $N$ \\
\midrule
1 & Head's occupation $\neq$ agriculture              & P4 & 12 \\
2 & \texttt{family\_type} solo / nuclear couple        & P2 & 11 \\
3 & Nuclear (or stable extended) with child in school  & P1 & 28 \\
4 & Remaining extended household                       & P3 & 14 \\
5 & Residual (joins the step-1 overrides)              & P4 & --- \\
\bottomrule
\end{tabular}
\end{table}

Two branches carry logic that the table cannot fully convey. The occupation override at step~1 takes precedence over all demographic criteria: any household not headed by an agricultural occupation routes to Profile~4 regardless of family structure, capturing the occupational actors of Infrastructure Mobility, Dual-household Strategy, and Hardware Resignation before any demographic test is applied. The educational branch at step~3 admits an extended-family household to Profile~1 only if it is \emph{stable}---no temporary migration and the solar home system kept at home---so that a settled, multi-generational household organised around school-aged children stays in Profile~1 rather than defaulting to Profile~3 on family-type code alone. Extended households that fail this stability clause, typically through active migration or a portable system, fall to Profile~3 at step~4.

Profile~4's population therefore combines two distinct mechanisms: households whose occupation structurally decouples them from the settlement's routine (step~1), and households whose family structure resolves cleanly into neither the stable-extended nor the solo/couple categories (step~5).\footnote{The count of 12 in Table~\ref{tab:tree} is the combined total of the step-1 and step-5 populations.}

Beyond the distribution shown in Table~\ref{tab:tree}, one household required the missing-data sensitivity check described in the methods, and zero households remained unclassified. Only a subset of each profile's survey-classified population was also interviewed, and only that subset supplies the qualitative material: \gls{p1} ($n_{\text{interview}}=18$, 64\% coverage), \gls{p2} ($n_{\text{interview}}=9$, 82\%), \gls{p3} ($n_{\text{interview}}=9$, 64\%), and \gls{p4} ($n_{\text{interview}}=6$, 50\%). The fourth profile's lower coverage is a specific limitation worth mentioning directly: it is also the profile carrying the dual-residence and structural-absence claims (Rule 10) that most need qualitative grounding, so its parametrization rests on the thinnest evidentiary base of the four.

Independently of the classification mechanism, correlational analysis of the resulting groups' quantitative survey data shows the profiles separating along behavioral outcomes that played no role in assignment, evidence consistent with convergent validity for the governing rules. Family type is strongly associated with daytime phone charging ($\eta = 0.76$ for a secondary device, $\eta = 0.72$ for a primary device), agricultural labor hours ($\eta = 0.67$), and cooking duration ($\eta = 0.56$), all large effects ($\eta^2 \approx 0.31$--$0.58$) consistent with \gls{p3}'s competitive USB-port and extended-cooking signature. Elderly presence is likewise associated with a delayed, flatter morning lighting curve ($\eta = 0.517$, $\eta^2 \approx 0.27$) and, more directly, with a lower count of working-age adults in the household ($r = -0.441$, $p < 0.001$), consistent with \gls{p2}'s demographic isolation. Given the sample size, these estimates should be read as indicative of large, robust effects rather than precise population values, and the working-age-adult correlation in particular reflects household composition as much as behavior.


\subsection{Virtual appliance parametrization}
\label{sec:va_parametrization}

For each \gls{ebp}, the qualitative rules governing its practices together with the quantitative data from the surveys are translated into \gls{va} parameters following the derivation protocol described in \ref{sec:va_parametrization}. Table \ref{tab:va-all} shows the appliance use parameters for all the profiles.

\begin{table}[H]
\centering
\footnotesize
\setlength{\tabcolsep}{4pt}
\caption{Virtual Appliance parameters across all four Energy Behavior Profiles.}
\label{tab:va-all}
\begin{tabular}{llcccccccc}
\toprule
\textbf{Profile} & \textbf{VA} & \textbf{Window} & \textbf{Rigidity} & \textbf{Power} & \textbf{occ.\_use} & \textbf{func\_time} & \textbf{func\_cycle} & \textbf{t\_f\_r\_v} & \textbf{rand\_var\_w} \\
\midrule
\multirow{9}{*}{1: Core}
& VA1 & 04:00--07:00 & Flexible & 3W & 0.35 & 84  & 42  & 0.2 & 0.30 \\
& VA2 & 07:00--18:00 & ---  & --- & --- & --- & ---  & --- & --- \\
& VA3 & 18:00--22:00 & Strict   & 3W & 1.00 & 192 & 130 & 0.10 & 0.20 \\
& VA4 & 22:00--04:00 & Strict   & 3W & 0.30 & 336 & 278 & 0.10 & 0.15 \\
& VA5 & 05:00--07:00 & Flexible & 2W & 0.20 & 84  & 42  & 0.20 & 0.30 \\
& VA6 & 07:00--18:00 & -- & -- & -- & -- & -- & -- & -- \\
& VA7 & 18:00--22:00 & Flexible   & 2W & 1 & 168 & 84  & 0.20 & 0.30 \\
& VA8 & 22:00--04:00 & -- & -- & -- & -- & -- & -- & -- \\
& VA9 & 00:00--24:00 & Chaos       & 2.5W & 1.00 & 480 & 120 & 0.20 & 0.00 \\
\midrule
\multirow{9}{*}{2: Elderly}
& VA1 & 04:00--07:00 & Chaos    & 3W & 0.15 & 117 & 54  & 0.30 & 0.35 \\
& VA2 & 07:00--18:00 & -- & -- & -- & -- & -- & -- & -- \\
& VA3 & 18:00--22:00 & Chaos    & 3W & 1 & 156 & 70  & 0.30 & 0.35 \\
& VA4 & 22:00--04:00 & Strict   & 3W & 0.15 & 288 & 240 & 0.10 & 0.15 \\
& VA5 & 04:00--07:00 & -- & -- & -- & -- & -- & -- & -- \\
& VA6 & 07:00--18:00 & -- & -- & -- & -- & -- & -- & -- \\
& VA7 & 18:00--22:00 & -- & -- & -- & -- & -- & -- & -- \\
& VA8 & 22:00--04:00 & -- & -- & -- & -- & -- & -- & --\\
& VA9 & 08:00--21:00 & Chaos    & 2.5W & 0.60 & 240 & 120 & 0.30 & 0.35 \\
\midrule
\multirow{9}{*}{3: Hub}
& VA1 & 04:00--07:00 & Flexible & 3W & 0.35 & 126 & 63  & 0.20 & 0.30 \\
& VA2 & 07:00--18:00 & -- & -- & -- & -- & -- & -- & -- \\
& VA3 & 18:00--22:00 & Flexible & 3W & 0.90 & 168 & 84  & 0.20 & 0.30 \\
& VA4 & 22:00--04:00 & Strict   & 3W & 0.50 & 288 & 240 & 0.10 & 0.15 \\
& VA5 & 04:00--07:00 & Strict & 2W & 0.35 & 126 & 63  & 0.10 & 0.20 \\
& VA6 & 07:00--18:00 & --       & -- & -- & --  & --  & -- & -- \\
& VA7 & 18:00--22:00 & Flexible & 2W & 1.00 & 168 & 84  & 0.20 & 0.30 \\
& VA8 & 22:00--04:00 & Strict   & 2W & 0.35 & 288 & 240 & 0.10 & 0.15 \\
& VA9 & 00:00--24:00 & Chaos    & 2.5W & 1.00 & 480 & 240 & 0.30 & 0.00 \\
\midrule
\multirow{9}{*}{4: Breakers}
& VA1 & 04:00--08:00 & Flexible & 3W & 0.35 & 90  & 55  & 0.20 & 0.30 \\
& VA2 & 08:00--18:00 & -- & -- & -- & -- & -- & -- & -- \\
& VA3 & 18:00--22:00 & Flexible   & 3W & 1.00 & 168 & 84  & 0.20 & 0.30 \\
& VA4 & 22:00--04:00 & Strict   & 3W & 0.40 & 288 & 240 & 0.10 & 0.15 \\
& VA5 & 04:00--07:00 & -- & -- & -- & -- & -- & -- & -- \\
& VA6 & 07:00--18:00 & -- & -- & -- & -- & -- & -- & -- \\
& VA7 & 18:00--22:00 & Flexible & 2W & 1.00 & 168  & 84  & 0.20 & 0.30 \\
& VA8 & 22:00--04:00 & -- & -- & -- & -- & -- & -- & -- \\
& VA9 & 00:00--24:00 & Chaos       & 2.5W & 1.00 & 240 & 120 & 0.20 & 0.00 \\
\bottomrule
\end{tabular}

\vspace{2pt}
{\footnotesize VA1--VA4: primary bulb, located in the main room used for several tasks. VA5--VA8: secondary bulb located outdoors or in the secondary room, same order. VA9: portable-device charging, including mobile phones, flashlights and radios. A dash indicates the VA is not used in the profile for diverse reasons described in the truth files. \texttt{t\_f\_r\_v} = \texttt{time\_fraction\_random\_variability}; \texttt{rand\_var\_w} = \texttt{random\_var\_w}. VA9 additionally carries \texttt{p\_var = 0.2} (declared default) in every profile, not shown here.}
\end{table}

The parameters for P1 are organized around one non-negotiable appliance: the indoor evening light, classified Strict, anchored firmly to dinner and homework with no reported drift in timing. Every other practice arranges itself around it --- the morning lights and outdoor evening light stay Flexible and sporadic, and the overnight light is Strict yet infrequent. Portable charging is Chaos-classified and highly frequent, spread across the whole day. Profile 1's signature is rigidity concentrated on one bulb rather than diffused across the schedule.

The second profile is defined by reduction to a single functional bulb: the secondary fixture is retired, since no lone- or two-person elder household reported running two lights at once. Its timing is the most erratic of any profile --- morning and evening light are both Chaos-classified, with the highest random variability in the study --- set against one island of rigidity: the overnight light, the least frequent practice here yet Strict, running unbroken from onset to dawn. Charging is the lowest and most time-bounded of all profiles, confined to daytime hours.

The multi-generational household, is the only profile to run both overnight lights at once. It also carries the highest overnight frequency and the longest charging cycle, reflecting the largest, most device-dense household, while evening and morning light stay Flexible to accommodate several co-resident generations. P3 is the profile of additive demand --- more lights, more devices, longer cycles --- rather than of any single rigid practice.

P4's defining feature is not encoded in any appliance. Its shapes resemble Profile 1 --- saturated evening light, chaotic charging --- but the evening light is Flexible rather than Strict, with lower functioning time under higher variability. The distinguishing behavior is carried instead by the household-level \texttt{prob\_home} mask (0.78, dropping to 0.61 in the planting and free-grazing seasons), which suppresses every appliance at once, halving charging time. The signature here is structural discontinuity at the household scale, not erraticism within any one practice.

Results reveal two patterns. First, \emph{rigidity tracks the practice, not the household type}. Evening light spans the full rigidity spectrum --- Strict in P1, Chaos in P2, Flexible in P3 and P4 --- so the same dinner/homework practice is classified differently depending on household structure. Overnight light, by contrast, is Strict in every profile that practices it at all: a household erratic about its evening light is nonetheless rigid about its overnight light once switched on.

Second, \texttt{occasional\_use} discriminates between profiles most sharply on the overnight light, ranging from 0.15 to 0.50 --- the widest profile separation of any lighting practice. Evening light clusters tightly (0.90--1.00 everywhere), so its frequency is far less diagnostic of profile identity than the qualitative rules would predict. Charging shows the widest spread of any appliance, from 0.60 to 1.00, tracking device ownership more directly than any single rule --- consistent with the Digital Delegation dynamic (Rule 5) operating across profiles rather than within one.

\subsection{Simulation and validation}
\label{sec:energy_demand_model}

The validation is reported in two stages, matching the two scales at which the models are actually used. First, the Socio-Technical Model is validated per profile against household measured data baselines; the Survey-Based Model enters here only as a scalar reference point, because its real failure mode is a community-scale phenomenon. Second, both models are scaled to the full classified population and compared as community aggregates, where the consequences of replacing a diverse population with a single replicated profile become visible.

The validation window is fixed to May, the month with continuous data across the logger set and the month falling between the community's two mobility peaks (the growing season, February--April, and the free-grazing/migration season, July--September), which minimizes the confounding effect of household absence on the load-curve comparison. The year-long traces of the settled (Profile~1) households confirm that May is representative for the schedule-bound majority, showing stable system use across all four agricultural seasons (Section \ref{sec:temporal_dynamics}), where the seasonal dynamics of the more mobile profiles are examined separately.

Results of the simulation and per-profile validation are shown in \ref{fig:load_curves}.

\begin{figure}[htbp]
    \centering
    \includegraphics[width=0.7\linewidth]{manuscript/Figures/Figure_4.png}
    \caption{Measured vs.\ simulated (Socio-Technical Model) mean daily load curves, by Energy Behavior Profile.}
    \label{fig:load_curves}
\end{figure}

Table~\ref{tab:perprofile_validation} summarizes the Socio-Technical Model's fit to household measured data across the four \glspl{ebp}, using a set of indicators covering time-series fidelity (nRMSE, LCSS, MPDADA), peak and energy magnitudes (P95, daily energy), and the day-to-day structure (MRSD, peak-time spread).

\begin{table}[H]
\centering
\caption{Per-profile validation of the Socio-Technical Model against household measured data, by Energy Behavior Profile. nRMSE is RMSE as a percentage of the profile's mean measured power; the value in parentheses is the household's own day-to-day coefficient of variation (CV), the natural-variability reference (nRMSE\,$\le$\,CV means the model's error is no larger than the household's own daily fluctuation). P95 and Daily-energy columns give the simulated-vs-measured relative error; MRSD and Peak-time spread are reported as measured\,/\,simulated, since it is the across-profile ordering, not the absolute value, that carries the socio-technical signal.}
\label{tab:perprofile_validation}
\begin{tabular}{lcccc}
\toprule
Metric & P1 & P2 & P3 & P4 \\
\midrule
nRMSE \% (CV \%)                    & 37.5 (23.2) & 44.2 (52.4) & 27.8 (30.0) & 37.6 (42.1) \\
LCSS                                & 0.875 & 0.969 & 0.865 & 0.927 \\
MPDADA                              & 0.238 & 0.004 & 0.048 & 0.435 \\
P95 power error                     & 3.8\% & 5.0\% & 5.5\% & 13.5\% \\
Daily energy error                  & 4.9\% & 6.7\% & 5.9\% & 0.4\% \\
MRSD (meas.\ / sim.)                & 0.232 / 0.180 & 0.463 / 0.574 & 0.300 / 0.260 & 0.421 / 0.469 \\
Peak-time spread, h (meas.\ / sim.) & 3.30 / 6.43 & 6.02 / 7.03 & 6.57 / 8.53 & 4.74 / 7.69 \\
\bottomrule
\end{tabular}
\end{table}

Overall error is low and interpretable against each household's own variability. The normalized RMSE ranges from 27.8\% (P3) to 44.2\% (P2), and for three of the four profiles it sits at or below the household's own coefficient of variation (Table~\ref{tab:perprofile_validation}): the model's typical deviation is no larger than the day-to-day fluctuation the real household already exhibits---the appropriate benchmark for an intermittent, appliance-level single-household signal, rather than the sub-20\% figures expected of smoothed, aggregated grid loads. Only the schedule-bound P1 exceeds its own CV, as expected: a rigid routine leaves little natural variability for error to hide within. Chronological shape is reproduced well throughout (LCSS 0.865--0.969), and peak sharpness (MPDADA) is best preserved for the low-consumption P2 and the multi-generational P3, loosest for the highly mobile P4 (Table~\ref{tab:perprofile_validation}).

The sizing-relevant quantities hold up across profiles. P95 power---the level exceeded only 5\% of the day, which drives inverter and peak-dispatch sizing---is reproduced to within 3.8--5.5\% for P1--P3 and 13.5\% for P4, a modest over-estimate. Aggregate daily energy is captured to within 0.4--6.7\% for every profile. The residual error concentrates in the daytime window rather than in the evening peak or aggregate totals, loosest for P1 and P3 (Table~\ref{tab:perprofile_validation}). One structural caveat not visible in the table is P3's absolute instantaneous maximum, which the model overshoots (7.06\,W measured vs.\ 10.88\,W simulated); because this single-hour spike leaves P3's P95, mean power, and daily energy essentially unaffected, it reads as a Virtual-Appliance modeling limitation rather than a distributional failure.

Regarding the day-to-day irregularity index (MRSD), the model reproduces the measured rank ordering of the four profiles---from the most schedule-bound to the most irregular---in both measured and simulated series (Table~\ref{tab:perprofile_validation}), under-predicting the amplitude slightly for the two rigid profiles and over-predicting it for the two irregular ones while preserving the rank. The peak-time spread shows the same signature: the schedule-bound P1 has the tightest measured peak timing and the model correctly places it lowest, though it over-disperses the spread for every profile.


\begin{figure}[htbp]
    \centering
    \includegraphics[width=0.7\linewidth]{manuscript/Figures/Figure_5.png}
    \caption{Load Duration Curve by profile: measured (solid black) vs.\
    Socio-Technical simulated (dashed navy) power, each sorted descending against
    the fraction of the day at or above that level. Discarding timing isolates how
    well the model reproduces the distribution of demand levels a household sits
    at. \textit{Provisional --- panels to be regenerated against this validation
    run before submission.}}
    \label{fig:ldc}
\end{figure}

Figure~\ref{fig:ldc} complements the chronological indicators with the \gls{ldc} for each profile: measured and simulated power sorted in descending order against the fraction of the day at or above that level. Because the sorted curve discards timing entirely, it isolates whether the model reproduces the overall \emph{distribution} of demand levels. The same \gls{rmse} statistic computed on this curve is markedly lower for every profile than the chronological \gls{rmse} of Table~\ref{tab:perprofile_validation}, so a substantial share of each profile's point-by-point error is attributable to \emph{when} demand occurs rather than to the distribution of demand levels the model reproduces. Normalized against mean power, the residual distributional error is small for all four profiles (11.7--16.9\%), confirming that the sorted-curve shape is well captured once timing is set aside; the sizing implications are developed in Section~\ref{sec:sizing}.

For reference, the same data scored against the single community-average profile of the Survey-Based Model yields nRMSE of 43--155\% across the four groups (against 28--44\%) and a structurally flat MRSD of 0.100 that cannot vary by household type. The per-household penalty of a non-adaptive profile is thus already visible; its full consequences emerge only once households are aggregated, to which we now turn.

Sizing decisions act on the community aggregate, not on individual households, and aggregation introduces an effect invisible at the single-profile scale and absent by construction from a homogeneous model: the partial cancellation of non-coincident peaks. Scaling both parametrizations to the classified population ($N=65$)---the Survey-Based Model as a single average profile replicated across all households (a \emph{homogeneous} community), the Socio-Technical Model as the four \glspl{ebp} in their surveyed proportions ($28{:}11{:}14{:}12$, a \emph{heterogeneous} community)---makes that effect measurable (Table~\ref{tab:community_validation}). Because the heterogeneous aggregate is the community-scale composition of the per-profile models already validated against real data (Table~\ref{tab:perprofile_validation}), it serves as the empirically anchored reference; community-scale metering is not available for an independent check, so this is an aggregation of validated models rather than a separate measured data validation.

\begin{table}[H]
\centering
\caption{Community-scale aggregation ($N=65$): the Survey-Based (homogeneous) Model against the Socio-Technical (heterogeneous) reference. Error is the homogeneous Model relative to the reference; sign indicates over- ($+$) or under-($-$) prediction. The coincidence factor is the aggregate peak divided by the sum of individual household peaks (mean\,$\pm$\,SD across days).}
\label{tab:community_validation}
\begin{tabular}{lrrr}
\toprule
Metric & Heterogeneous (ref.) & Homogeneous (SB) & Error \\
\midrule
Coincident peak (W)            & 311.31 & 345.05 & $+10.8\%$ \\
Coincidence factor             & $0.717 \pm 0.056$ & $0.774 \pm 0.030$ & $+8.0\%$ \\
P95 power (W)                  & 207.86 & 260.50 & $+25.3\%$ \\
Daily energy (Wh)              & 1849.4 & 2233.4 & $+20.8\%$ \\
Base / night load, 00--04 (W)  & 72.69  & 39.89  & $-45.1\%$ \\
Morning load, 04--10 (W)       & 58.06  & 98.88  & $+70.3\%$ \\
MRSD (chaos index)             & 0.150  & 0.016  & $-89.3\%$ \\
\bottomrule
\end{tabular}
\end{table}


The coincidence factor makes the diversity effect explicit. In the heterogeneous reference it averages 0.717: individual peaks land at staggered times and only partially stack, so the aggregate peak sits well below the sum of household peaks. The homogeneous model raises this to 0.774 and, tellingly, makes it more than half as variable from day to day ($\pm0.030$ vs.\ $\pm0.056$), a first sign of the variability collapse that recurs throughout. The inflated coincidence carries straight into the coincident peak (345\,W vs.\ 311\,W) and the P95 level ($+25.3\%$): a system sized on the homogeneous curve would over-provision peak capacity by roughly a quarter.

Beyond the peak, the daily-shape errors run in \emph{opposite} directions---the aggregate signature of replicating one profile in place of a diverse population. Because identical households wake synchronously, the homogeneous aggregate produces a phantom dawn spike near 06:00 (Figure~\ref{fig:community_shadows}) that overstates the morning window by 70.3\%; in the reference, staggered morning anchors across profiles---the early waking of the Agricultural Core (Rule~2) offset against the later, flatter waking of the Isolated Elder---smear that spike into a lower, broader shoulder. Conversely, the homogeneous aggregate is base-load blind: it under-predicts overnight load by 45.1\%, because a single average profile omits the safety lighting that the overnight-lighting profiles contribute. Diversity fills troughs and shaves peaks; replicating one profile exaggerates both, accumulating over the day into a $20.8\%$ over-prediction of total energy---i.e.\ systematic over-provisioning.

\begin{figure}[htbp]
\centering
\includegraphics[width=0.8\linewidth]{manuscript/Figures/Figure_6.png}
\caption{Daily-trace envelopes.}
\label{fig:community_shadows}
\caption{(a) Community-aggregate Load Duration Curve, heterogeneous reference vs.\
homogeneous Model. (b) Daily-trace envelopes: the heterogeneous reference spreads
across a wide day-to-day band, the homogeneous Model a narrow, near-deterministic
one (MRSD 0.150 vs.\ 0.016). Provisional styling.}
\label{fig:community_variability}
\end{figure}

Finally, day-to-day variability collapses: the homogeneous aggregate's MRSD is 0.016 against the reference's 0.150, an 89.3\% reduction (Figure~\ref{fig:community_variability}). A single-profile community produces near-identical days, with none of the structural irregularity that a mixed population injects into the aggregate. This is the community-scale counterpart of the flat, by-construction MRSD seen per household: a homogeneous model has no mechanism to express behavioral irregularity at either scale. These distortions are not scattered inaccuracies but the specific, physically interpretable consequences of replacing a diverse population with one replicated profile, and they carry directly into system sizing.

\section{Discussion}
\label{sec:discussion}

The results support the claim that qualitative analysis produced behavioral parameters the survey data could hardly have generated, and those parameters measurably reshaped the simulated load curves. The clearest case is the rigidity classification. The Survey-Based Model takes a common global approach, returning a single irregularity value for every household , whereas the Socio-Technical Model reproduces the real cross-profile ordering, capturing heterogeneity better. The claim is not that qualitative methods are universally necessary, but that for micro-level behavioral parameters in data-scarce, culturally-specific off-grid settings, qualitative understanding is the input without which the model cannot be built.

\subsection{Temporal dynamics}
\label{sec:temporal_dynamics}

The validation in Section~\ref{sec:energy_demand_model} uses measured data from May. A single-month comparison is bounded in two ways, which the extended-coverage dataloggers and the revisited households let us address in turn: whether it captures how differently households behave at the moment of measurement, and whether that behavior is stable over the system's lifetime. The first is heterogeneity across households; the second, drift within them.

The extended traces resolve the cross-sectional contrast directly. The \gls{p1} households with long coverage use the system every recorded day, with no multi-day interruption at any point in the year; their annual traces resemble one another and their daily energy varies only moderately around a consistent level. This regularity is the physical counterpart of the profile's governing rules. The \gls{p4} present the opposite condition, but not uniformly: one household is absent for a large fraction of its valid days across recurring multi-day interruptions, another is continuously present at a persistently low level, a third is steady throughout. What distinguishes P4 is therefore not that every household is erratic but that the profile has \emph{no shared load signature}---its members are grouped by the social conditions that decouple them from the settlement's routine (labor migration, dual-household residence, occupational mobility; Rules~8, 10, 13), and those conditions express differently in each. Predictability, not a specific pattern, is what P1 has and P4 lacks. The consistent finding is a division between the settled Profile~1, whose households are mutually similar and individually regular, and the mobile and extended profiles, whose households are mutually dissimilar.

This annual-scale division is independently corroborated within the validation month itself: the Chaos Index is lowest for P1 and high for P4
(Table~\ref{tab:perprofile_validation}). The Socio-Technical Model reproduces this ordering; the Survey-Based Model returns a single value and cannot express it. Two independent routes---a full year of extended traces and a single validation month---arrive at the same structure, which is what justifies May as the validation month: it falls between the two periods of greatest mobility \citep{gaioc_del_territorio_de_raqaypampa_plan_2018}, and the settled Profile~1 households that dominate the validation set behave in May as they do year-round. The community-scale coincidence factor confirms this---$0.756$ in May against a $0.761$ settled-month average and $0.672$ across the mobility seasons. What the validation does not capture is the seasonal and aperiodic absence of the mobile profiles; that is a property of the annual scale alone.

Heterogeneity at a fixed moment is only one boundary; whether it is stable over time is another. The revisited households indicate that it is not, and that they drift along the same axis the profiles are built on, through two mechanisms. The first is hardware attrition: a radio reported non-functional at one visit remained unrepaired a year later, its use permanently contracted [P90]---the temporal form of the Hardware Resignation rule (Rule~13). But attrition is not deterministic; elsewhere a broken component was replaced or restored [P76], recovering the load. Whether a fault ends in resignation or restoration depends on household resources, so it cannot be modeled as a uniform decay rate. The second is demographic: as children age and adult children out-migrate, a household loses the members that define its profile and comes to satisfy another's criteria, drifting from the schedule-bound Core toward the flatter, elder-headed Isolated Elderly pattern. The settled and mobile poles are thus not fixed addresses but positions households move between over a system's lifetime.

Together these boundaries delimit the model's validity in composition and in time. The parametrization captures the community accurately at the moment of measurement, but a profile fixed at deployment will diverge over a ten-to-twenty-year lifetime as hardware degrades unevenly and households transit between profiles---representing which would require stochastic degradation terms admitting both resignation and restoration, and profile-transition probabilities \citep{XXX}. These are directions for future modeling, not corrections to a method whose scope is explicitly a validated snapshot. The extended evidence rests on seven households and the temporal evidence on five, offered as descriptive complements to the validation rather than independent statistical results (Section~\ref{sec:limitations}).

\subsection{Implications for off-grid system sizing}
\label{sec:sizing}

Off-grid sizing is determined by extreme conditions rather than average ones, which is why a model that reproduces mean demand can still yield a mis-sized system \citep{sizing-consequences}. Four demand properties govern the sizing of a solar system: daily energy, which sets PV and battery capacity; peak power and its duration, which set the converter rating; day-to-day variability, which sets the required storage autonomy; and the temporal alignment between demand and generation, which sets how much demand must be met from storage \citep{sizing-methodology}. Each is resolved per profile by the Socio-Technical Model and collapsed to a single value by the Survey-Based Model.

Variability illustrates the distinction. Represented by one averaged profile, the Survey-Based Model provides no basis for distinguishing a household with stable demand from one whose demand is intermittent. The Socio-Technical Model instead reproduces the measured ordering of irregularity across profiles, erring in both directions around it---itself evidence that the model tracks real household heterogeneity rather than applying a uniform correction factor. Because variability determines the storage autonomy a system requires, this ordering is a sizing input, and it originates in what the qualitative approach captured.

Peak magnitude follows a more graded pattern. The Survey-Based Model's fixed peak approximates the measured value closely for P1, is moderately off for P3, and departs sharply for P2 and P4, because a single value can coincide only with households near the mean (Table~\ref{tab:perprofile_validation}). The Socio-Technical Model holds curve-peak error within $0.1$--$14.1\%$ across all four profiles. The P95 power level---the demand exceeded $5\%$ of the day, and a direct basis for converter sizing---shows the same asymmetry: the Survey-Based Model returns $4.06$~W for every profile, matching the measured value for P1 and P3 ($4.23$ and $3.90$~W) while over- or under-sizing P2 and P4 (measured $2.28$ and $2.88$~W; Table~\ref{tab:ldc}). A system dimensioned from the averaged model is adequate for households resembling the community mean and mis-sized for those that do not.

Peak \emph{timing} behaves differently. The two models are closely comparable on modal peak hour, with no offset exceeding one hour for any profile, because an evening-anchored peak is common to all four and a single shared evening window places it in the correct interval regardless of parametrization. The residual one-hour offset the Socio-Technical Model shows for P2--P4 is partly mechanistic: the model represents each evening block with a single symmetric variability parameter, whereas the telemetry shows a synchronized onset followed by a more dispersed offset, and an over-gradual symmetric decay shifts the simulated peak later (Section~\ref{sec:limitations}). This has limited consequence for individual-SHS sizing, which is set by peak magnitude rather than exact hour.

That consequence changes with aggregation. For an individual SHS there is no aggregation: each system meets its own household's peak and variability, so modeling error propagates directly into its hardware. For a microgrid, the simultaneous peak of connected households is lower than the sum of individual peaks, and estimating this coincidence factor requires exactly the profile-level differences in peak timing and shape that a single averaged profile cannot represent \citep{XXXX}. The peak-timing detail of limited consequence for individual SHS therefore becomes a determinant of microgrid
sizing.

Aggregating both models to the community scale quantifies this. Over a simulated year the heterogeneous community averages a coincidence factor of $0.72$ (range $0.56$--$0.83$, $\sigma=0.056$) against $0.77$ for the homogeneous one (range $0.69$--$0.88$, $\sigma=0.030$), with a coincident peak of $311$~W versus $345$~W (Table~\ref{tab:community_validation}). Identical households peak in near-unison, so the homogeneous model both inflates the summed individual peaks and cancels less of them, overstating simultaneity twice over and, for a microgrid, sizing shared transformer and generation capacity for a peak the real population never reaches: staggered peaks across profiles reduce the aggregate, and a model that collapses them into one shape cannot reproduce that reduction. The gap is also seasonal in a way the homogeneous model cannot express---the heterogeneous factor falls to $0.672$ across the mobility seasons and recovers to $0.761$ when the community is settled, while the homogeneous factor stays flat year-round (Section~\ref{sec:temporal_dynamics}).

The community Load Duration Curve (Fig.~\ref{fig:community_shadows}) captures both errors in one view: the homogeneous curve sits above the reference through the upper fraction of the day and below it in the sustained-load tail, crossing near the 80th percentile. A system sized from it would over-provision PV and converter capacity while under-sizing storage, whose overnight base load it under-predicts by $45.1\%$. This second error is decisive: base load sets required autonomy, and
it stems entirely from the overnight safety-lighting behaviour (Rule~7) the homogeneous profile omits. The homogeneous model thus errs in \emph{opposite directions} on the two quantities that govern off-grid sizing---an effect no recalibration of a single average profile can fix, because it stems from missing diversity rather than a mis-set average.

\subsection{Reading low demand correctly}
\label{sec:reading_low_demand}

A load curve that is low or flat admits more than one explanation, and the three the Raqaypampa evidence shows imply different, sometimes opposite, planning responses. Within the Energy Cultures Framework, material culture shapes the range of practices a household can perform \citep{stephenson_energy_2015}; the evidence here specifies that further. A practice can be absent from the evening load because the technology available for it is inadequate for the task, because the demand it would generate is being met elsewhere, or because the household has arrived at a genuine sense of sufficiency and does not seek more.

The range of afforded practices is governed not only by the presence of a technology but by the adequacy of the service it provides. Evening textile work is the clearest instance: across several households, adult women reported that weaving and spinning are done in daylight rather than under the SHS light, because the illumination is insufficient for detailed handwork---``I do not do any activity with the light because it is not clear enough'' [P72]; ``during the day we weave; we do not use the light'' [P74] (Rule~14). Provision of a basic bulb is therefore not sufficient to generate night-time productive activity; the quality of the light determines whether the practice is afforded at all. The load curves parametrized here describe demand as constrained by prevailing material conditions, not the latent demand that might arise were they altered---so the absent evening productive-lighting load is not a modelling gap awaiting correction but a correctly represented consequence of an inadequate affordance (Section~\ref{sec:va-parametrization}). A model inferring latent demand from device ownership alone would size for a weaving load the material conditions do not permit.

Mobility and multi-locality are not anomalies in remote rural contexts but established livelihood strategies: seasonal labour migration and the maintenance of more than one residence are widespread responses to dispersed land, uneven infrastructure, and precarious incomes across the rural Global South
\citep{rural-multilocal-livelihoods, energy-stacking}. A demand model that assumes each household occupies a single fixed dwelling therefore encodes an assumption that does not hold for a meaningful share of the population, and the consequence is not a small numerical error but a systematic misreading of what a load curve means. The clearest form here is dual-residence, documented with dated corroboration in five households (Rule~10): one maintains two separate systems and relocates entirely to a grid-connected second dwelling for a defined part of the year, leaving the \gls{shs} unoccupied [P76]. The mechanism is legible in the households' own accounts---higher-consumption practices are carried out ``at our other house \ldots{} where there is ELFEC [grid] light,'' while the SHS dwelling is used only ``to cook dinner, eat, and then sleep'' [P54]. There, such a household registers a low, flat curve that a conventional analysis reads as low demand or an adequately-sized system; the interview evidence shows the opposite. The demand exists and is substantial, but is metered at the second residence. The flatness is not the absence of demand but its spatial displacement.

A third condition is distinct from both: several households frame present consumption against remembered pre-electrification conditions---the kerosene lamp (\emph{mechero}) and its smoke---and describe current service as adequate rather than deficient (Rule~16). This energy-poverty-legacy frame, common among older respondents, is associated with a conservative use of the resource, the light treated as a valued and finite good. The evidence is qualitative and is not advanced as a quantified determinant of demand; it indicates only that a low load may, in some households, reflect a settled sufficiency rather than a shortfall a larger system would fill.

The three conditions are diagnosed differently but converge on one requirement: occupancy and movement cannot be read from a single survey field. The same flat curve is produced by a genuinely low-consuming resident household, a dual-residence household absent for part of the year, and a household with an inadequate but present affordance---and the migration survey field did not discriminate them, being common across the sample without identifying genuine dual-residence (Section~\ref{sec:ebp-classification}). Yet the three pull system design in different directions. A dual-residence household's suppressed load is latent and recoverable through capacity: a grid connection or larger system at the SHS site could draw back the demand met elsewhere, a step increase curve-fitting on present consumption cannot anticipate. An inadequate-affordance household's load is also latent but recoverable through a different lever entirely---quality of provision rather than added capacity, so a latent-demand reading that simply adds generation would size for the wrong intervention \citep{azimoh_illuminated_2015}. A genuinely sufficient household's low load is not latent at all, and sizing for a shortfall it does not perceive over-provisions without cause. Reading a flat curve correctly---identifying which condition produced it---is therefore a prerequisite for sizing correctly, and it requires the qualitative account of household structure that the load curve and the standard survey instrument do not supply.

\subsection{Scalability and limitations of the socio-technical pipeline}
\label{sec:scalability}

The main contribution of this work is a procedure as much as a result specific to Raqaypampa, and the two must be separated when considering where the approach transfers and where, even there, it falls short. What transfers is the extraction and translation method: the derivation protocol, the logic of the two crosswalks mapping respondent language to the probability and rigidity of activation, the triangulation of interview narrative against survey time-use and physical anchors, and the logic by which a qualitative typology becomes a deterministic classification of the entire surveyed population. What does not transfer is the content this method produced here---the specific behavioral rules, the four Energy Behavior Profiles, and the parameter values are particular to this community and carry no claim of generality beyond it.

Transferability is a property of the procedure, not of the crosswalk this study produced. The bins were developed from respondents' own words and each validated against quotes from multiple profiles rather than one (Section~\ref{sec:va-parametrization}; Supplementary~S2), which is what licensed collapsing four profile-specific crosswalks into one: a given expression carried approximately the same meaning across demographically distinct households. That cross-profile stability does not guarantee the same word-to-value mapping in a different linguistic setting, where frequency and timing language may calibrate
differently. The transferable object is therefore the derivation procedure---parameter-driven extraction, inductive binning, cross-profile validation, and iterative refinement against the transcript---which a new site re-runs on its own corpus to derive and validate its own crosswalk. What makes this reusable in practice is that its inputs are the ones available in data-scarce settings: interview transcripts and a basic survey, with no requirement for the metered consumption data typically absent where demand estimation is hardest \citep{data-scarcity}.

Because each parameter is tied to an explicit, declared source, the truth-file structure has prescriptive value: it specifies which observations a parameter requires, letting a future study design its instruments around that data rather than recovering it by triangulation after the fact, as was necessary here. Reuse
is nonetheless a re-derivation---it requires the new setting's own fieldwork, its own rule extraction, and its own validation of the bins against local language. The claim is accordingly bounded: the pipeline suits comparable data-scarce, culturally-specific off-grid contexts where interview and survey data are
available but telemetry is not, and is not advanced beyond that class of setting.

The limitations cluster in three areas: the evidentiary base, the simulation engine, and the scope of what was asked. The 100-household universe is unevenly covered by qualitative material. Parameters are derived from each profile's interviewed subset and applied uniformly to the rest, so a profile's reliability
tracks its coverage---thinnest for the System Breakers (50\%, six of twelve households), the profile whose defining dual-residence claims most depend on qualitative grounding and whose parameters therefore rest partly on pooled rather than per-household evidence (Supplementary~S1). A few classifications rest on a documented analyst override where qualitative evidence contradicted the survey fields (Supplementary~S2), evidencing the say-do gap. The evidence for annual dynamics is thinner still and is offered as descriptive, not statistical. Seasonality is the sharpest instance: the data locate the dominant seasonal driver in occupancy and generation rather than appliance behavior, and structural absence is modeled by a household-level occupancy mask---but the anchoring absences are too few and too coarsely dated to resolve a household's contiguous multi-month spell rather than its average absence rate. This round makes the missing measurement precise rather than settling the seasonal profile; future fieldwork should record absence durations and their timing per household.

Two limitations sit in RAMP, and they are linked. The most consequential emerged from validation: evening lighting's onset is tightly consistent across households and days, while its offset varies far more. The model cannot reproduce this asymmetry, because its variability parameter governs each activation block as a
single symmetric unit. This is also a ceiling on self-report---``we use the light from six to ten'' gives a modal window, not a distribution over nights \citep{sensor-vs-survey}---and its measured cost is only the residual one-hour peak-timing offset on three of four profiles (Section~\ref{sec:sizing}). Representing it would mean splitting one appliance into two Virtual Appliances with independent onset and offset variability, which runs into the second limitation: RAMP provides no mutual-exclusion mechanism, so two Virtual Appliances sharing one device can overlap and double-count its load. That overlap is bounded, computable, and small enough not to affect the present results, but a native exclusivity (or asymmetric-window) mechanism would resolve both at once and let fieldwork probing onset-versus-offset variability feed straight into the model \citep{paper about RAMP}.

Finally, a boundary of what was asked rather than of method. Because the systems here are very small, the most consequential dynamics may lie beyond consumption: how a household builds the capacity to operate, maintain, and extend an unfamiliar technology, and how that knowledge is transferred, is a question this snapshot raises but does not answer. A longitudinal design pairing load measurement with indicators of adaptive capacity and knowledge transfer is the natural next study, and would also supply the repeat-absence and timing data the occupancy and seasonal extensions require.


\section{Conclusion}
\label{sec:conslusion}

This study set out to characterize how socio-cultural factors shape energy behavior in an Indigenous territory and to operationalize those insights for bottom-up demand modeling.

Applying the Energy Cultures Framework to the case of Raqaypampa, the study extracts seventeen socio-technical rules governing daily energy use and consolidates them into four Energy Behavior Profiles. Read together, these findings show that electricity demand here is shaped less by appliance ownership than by a cultural logic anchored in agricultural labor. Because the communities lack irrigation, survival is tethered to the rainy season, and this productive calendar sets the rhythm of energy use: the Agricultural Dictate pushes routines into the pre-dawn hours to capture daylight in the fields, the protected evening window sustains children's study after the working day, and seasonal migration and the relocation of panels to agricultural camps drive the extended absences that flatten some load curves entirely. Overlaying this seasonal rhythm is a slower trend—permanent rural-to-urban migration—that gradually reshapes household composition and, with it, the demand a system must serve. Energy practice is thus woven into the work of producing food and sustaining the household.

The four profiles are best understood not as fixed categories but as positions in a household life cycle. As children reach school age and adults migrate, a household drifts along the same demographic axis on which the profiles are built—typically from the schedule-bound Core toward the flatter, elder-headed pattern. The profiles therefore capture both the community's heterogeneity at a given moment and the trajectory a household follows over a system's lifetime.

Operationalizing these insights, the study integrates qualitative and quantitative data into the parameters of a stochastic demand model, disaggregating identical hardware into behavior-specific Virtual Appliances. Comparing this socio-technical model against a survey-only baseline shows that qualitative parametrization reshapes the simulated load curve in ways survey data alone cannot—reproducing the cross-profile ordering of behavioral irregularity, the chronological shape of daily routines, and aggregate daily energy that a single community-average profile cannot vary by household type. The comparison is deliberately bounded: the socio-technical model is not superior on every metric, but it is more informative for the properties on which system sizing and resilient design depend—peak magnitude, day-to-day variability, and displaced demand. In these data-scarce, culturally specific settings, then, qualitative understanding does not only supplement the quantitative model; it structures it—a single rigidity judgment, for instance, fixes three separate parameters that survey data could not have supplied. A stochastic engine such as RAMP is only as good as the care taken in its parametrization.

What extends beyond Raqaypampa is this procedure rather than its outputs. The four profiles and their parameter values are specific to the community, but the extraction-and-translation method requires interview transcripts and a basic socioeconomic survey, the inputs available precisely where metered data is not. Applying it elsewhere is a re-derivation rather than a redeployment, demanding local fieldwork and local validation of the crosswalk bins against local language.

The boundaries of the work point to three directions. First, the parametrization is a validated snapshot: it fixes each household in one profile, whereas hardware degrades unevenly and households transit between profiles over a system's twenty-year lifetime—capturing this would require stochastic degradation terms and profile-transition probabilities. Second, the evidence locates the dominant seasonal driver in occupancy rather than appliance behavior, but the anchoring absences are too few and too coarsely dated to resolve contiguous multi-month absences, so future fieldwork should record absence durations and timing per household. Third and most broadly, because these systems are very small, the most consequential dynamics may lie beyond consumption—in how households build and pass on the capacity to operate and maintain an unfamiliar technology. A longitudinal design pairing load measurement with indicators of adaptive capacity would address this directly, while supplying the absence and drift data the other two extensions require.

Within these bounds, the method offers a reproducible path toward demand models grounded in how communities actually live rather than how top-down assumptions presume they should—a prerequisite for electrification planning that is both technically sound and context-appropriate.

\section*{Declaration of competing interest}
The authors declare that they have no known competing financial interests or personal relationships that could have appeared to
influence the work reported in this paper

\section*{Acknowledgments}
The authors express their sincere gratitude to the Academy for Research and Higher Education (ARES, by its French acronym) from Belgium for their invaluable financial support in the framework of the Research for Development Project: \textit{Tailored Energy System Models for Energy Planning in Bolivia}. This project would not have been possible without their commitment to advancing research that addresses the unique energy challenges faced by developing countries.

\section*{Data availability}

The dataset supporting the findings of this study is openly available on...

\section*{Declaration of generative AI and AI-assisted technologies in the manuscript preparation process}
During the preparation of this work the author(s) used Claude (Anthropic) in order to improve the readability and language of the manuscript. After using this tool/service, the authors reviewed and edited the content as needed and take full responsibility for the content of the published article.

%% The Appendices part is started with the command \appendix;
%% appendix sections are then done as normal sections

\appendix
\section*{Appendix A. Supplementary material}

Supplementary material to this article can be found online at...



%% If you have bibdatabase file and want bibtex to generate the
%% bibitems, please use
%%
 \bibliographystyle{elsarticle-num} 
 \bibliography{cas-refs}

%% else use the following coding to input the bibitems directly in the
%% TeX file.

% \begin{thebibliography}{00}

% %% \bibitem{label}
% %% Text of bibliographic item

% \bibitem{}

% \end{thebibliography}
\end{document}
\endinput
%%
%% End of file `elsarticle-template-num.tex'.
